% year: תשפ"ד
% semester: ב
% moed: א
% number: 3
% categories: func-analysis
% source: exams/מבחנים/תשפד/מועד א תשפד סמסטר ב.pdf

\begin{question}
חקרו את הפונקציה הבאה (נק' קיצון מקומיות, תחומי קמירות/קעירות, אסימפטוטות):
\[ f(x)=\frac{2-x-3x^2}{x^2-1} \]
שרטטו את הגרף שלה.
\end{question}

\begin{hint}
פרקו את המונה והמכנה לגורמים — יש גורם משותף.
\end{hint}

\begin{hint}
אחרי הצמצום ניתן לכתוב $f(x)=-3-\frac{1}{x-1}$ (עבור $x\neq\pm1$).
\end{hint}

\begin{solution}
\textbf{תחום הגדרה:} $x^2-1\neq0$, כלומר $x\neq\pm1$.

\textbf{פישוט:} $2-x-3x^2=-(3x^2+x-2)=-(3x-2)(x+1)$ ו-$x^2-1=(x-1)(x+1)$. לכן לכל $x\neq\pm1$:
\[ f(x)=\frac{-(3x-2)(x+1)}{(x-1)(x+1)}=\frac{2-3x}{x-1}=-3-\frac{1}{x-1}. \]
(בדיקה: $-3(x-1)-1=2-3x$.)

\textbf{הנקודה $x=-1$:} $\lim_{x\to-1}f(x)=\frac{2+3}{-2}=-\frac52$ — גבול סופי, ולכן זו נקודת אי-רציפות סליקה ("חור" בגרף בנקודה $(-1,-\frac52)$), ואין בה אסימפטוטה אנכית.

\textbf{חיתוך עם הצירים:} $f(0)=-2$; $f(x)=0\iff x=\frac23$. הנקודות $(0,-2)$, $(\frac23,0)$.

\textbf{נגזרת ראשונה:} $f'(x)=\frac{1}{(x-1)^2}>0$ לכל $x$ בתחום. לכן $f$ עולה בכל אחד מהקטעים $(-\infty,-1)$, $(-1,1)$, $(1,\infty)$, ו\textbf{אין נקודות קיצון מקומיות} (הנגזרת לא מתאפסת בשום מקום, ולפי משפט פרמה אין קיצון פנימי).

\textbf{נגזרת שנייה:} $f''(x)=-\frac{2}{(x-1)^3}$.
\begin{itemize}
  \item עבור $x<1$ ($x\neq-1$): $(x-1)^3<0$, לכן $f''>0$ — $f$ \textbf{קמורה} ($\cup$) ב-$(-\infty,-1)$ וב-$(-1,1)$.
  \item עבור $x>1$: $f''<0$ — $f$ \textbf{קעורה} ($\cap$) ב-$(1,\infty)$.
\end{itemize}
אין נקודות פיתול: שינוי הקמירות קורה רק ב-$x=1$, שאינה בתחום ההגדרה.

\textbf{אסימפטוטות:}
\begin{itemize}
  \item אנכית: $x=1$. $\lim_{x\to1^-}f(x)=-3-\frac{1}{0^-}=+\infty$, $\lim_{x\to1^+}f(x)=-\infty$.
  \item אופקית: $\lim_{x\to\pm\infty}f(x)=-3$ (המעלות במונה ובמכנה שוות, יחס המקדמים המובילים $\frac{-3}{1}$). לכן $y=-3$ אסימפטוטה אופקית בשני הכיוונים; אין אסימפטוטה משופעת.
\end{itemize}

\textbf{הגרף:} היפרבולה $y=-3-\frac1{x-1}$ (הזזה של $-\frac1x$) עם חור בנקודה $(-1,-\frac52)$. משמאל ל-$x=1$ הגרף עולה מהאסימפטוטה $y=-3$ (מעליה) דרך החור $(-1,-2.5)$, $(0,-2)$, $(\frac23,0)$ ושואף ל-$+\infty$ כאשר $x\to1^-$, והוא קמור. מימין ל-$x=1$ הגרף עולה מ-$-\infty$ ומתקרב מלמטה ל-$y=-3$, והוא קעור.
\end{solution}
